Для оценок скорости методов одной выпуклости мало: нужны оценки функции параболами снизу и сверху. Сильная выпуклость (параметр $\mu$) дает нижнюю параболу и \hyperref[note:51CC880D-BF19-43E2-9599-FB31A9A7367F]{единственность минимума}, липшицевость градиента (параметр $L$) --- верхнюю. Для $f \in C^2$ оба условия записываются через гессиан: $\mu I \le \nabla^2 f(x) \le L I$, а отношение $\kappa = L/\mu \ge 1$ определяет скорость градиентного спуска.
